% =========================================================================
%  Congruence certificates for Collatz-like machines: a census, a sieve
%  completeness theorem, and the impossibility of separating two flagships
%
%  DRAFT v2 -- restructured July 31, 2026.
%  Design commitments of this draft:
%    * research goals stated in Section 1, progress assessed against them
%      in Section 1 (summary) and Section 9 (in full);
%    * every result carries an explicit epistemic label (Section 1.5) and
%      appears in the status ledger of Section 8;
%    * the distinction between Turing machines, counter machines and
%      Collatz-like functions is made precise in Section 1.1 and respected
%      throughout;
%    * claims are stated at the strength actually established, and the
%      limitations section (1.4) and threats to validity (8.2) are not
%      optional reading.
%
%  Compile: tectonic main.tex   (bibliography is inlined; see README)
% =========================================================================
\documentclass[11pt]{article}

\usepackage[margin=1.1in]{geometry}
\usepackage{amsmath,amssymb,amsthm}
\usepackage{booktabs}
\usepackage{xcolor}
\definecolor{linkblue}{RGB}{26,60,110}
\usepackage[colorlinks=true,linkcolor=linkblue,citecolor=linkblue,urlcolor=linkblue]{hyperref}

\newtheorem{theorem}{Theorem}[section]
\newtheorem{lemma}[theorem]{Lemma}
\newtheorem{corollary}[theorem]{Corollary}
\newtheorem{proposition}[theorem]{Proposition}
\theoremstyle{definition}
\newtheorem{definition}[theorem]{Definition}
\theoremstyle{remark}
\newtheorem{remark}[theorem]{Remark}
\newtheorem{question}[theorem]{Question}
\newtheorem{observation}[theorem]{Observation}

\newenvironment{evidence}
  {\par\smallskip\noindent\footnotesize\textbf{Evidence.}}
  {\par\smallskip}

\newcommand{\Z}{\mathbb{Z}}
\newcommand{\N}{\mathbb{N}}
\newcommand{\ord}{\operatorname{ord}}
\newcommand{\oddpart}{\operatorname{oddPart}}
\newcommand{\vp}[1]{v_{#1}}
\newcommand{\Proved}{\textsc{[proved]}}
\newcommand{\CAP}{\textsc{[proved, computer-assisted]}}
\newcommand{\MV}{\textsc{[machine-verified]}}
\newcommand{\Meas}{\textsc{[measured]}}
\newcommand{\Interp}{\textsc{[interpretation]}}

\title{Congruence certificates for Collatz-like machines:\\
a census, a sieve completeness theorem, and the\\
impossibility of separating two flagship orbits}
\author{Tom Zahavy\thanks{This research program was carried out with an AI
assistant (Claude, Anthropic).  All proofs were additionally machine-checked
as described in Section~\ref{sec:conventions}; the form of credit and the
description of the workflow are to be settled before submission.}}
\date{July 2026 --- working draft (v2)}

\begin{document}
\maketitle

\begin{abstract}
We study \emph{one-schema valuation machines}: Collatz-like functions on the
positive integers determined by five integers
$(\alpha,\beta,\gamma,\delta,\varepsilon)$, where for $x=2^{v}u$ with $u$ odd
the function halts if $u=1$ and otherwise, writing $u=2k+1$, returns
$A_v k+B_v$ with $A_v=\alpha 2^{v+1}+\beta$ and
$B_v=\gamma 2^{v}+\delta v+\varepsilon$.  The class is of interest because the
published halting-equivalent reduction of the \emph{Space Needle}---a $6$-state
Turing machine and a well-known open problem in the busy-beaver
community---is the member $(1,3,1,1,0)$.

We report four things.  (i) A census of the $1{,}077$ well-defined members of
a parameter box, of which $318$ are decided by an exact congruence procedure
that is complete for its certificate class, with every certificate
independently re-audited; and closed-form criteria for separation at moduli
$3,4,6,8$, each verified exhaustively on $(\Z_m)^5$.  (ii) A closed form for
the last-step sieve of the family, from which ten members are proved never to
halt after their first step, together with a \emph{completeness theorem}:
within the box those ten are all, every other applicable member carrying an
explicit positive witness that the method fails on it.  (iii) A chain of
saturation theorems ending in: if $\beta$ is odd, $A_v>0$ and
$\gcd(\delta,M')=1$ ($M'$ the odd part of $M$), then the congruence closure of
every residue is all of $\Z_M$, at every modulus.  Consequently \emph{no
congruence certificate at any modulus separates the Space Needle's reduction};
the same argument transfers to base $3$ and closes the corresponding question
for a second machine studied here.  (iv) A test of the theory outside its design set: a
second cryptid found independently by the busy-beaver community (the ``sheep''
machine, April 2026) lies inside the interface, its last-step sieve closes
completely---confirming a prediction made beforehand about which machines it
should close for---and its halting set is obtained in closed form.  (v)
Expressiveness barriers for the class, and a refutation of counting-style
arguments against its universality.

We are explicit about what is \emph{not} achieved.  No Turing machine's
halting problem is decided here; neither flagship orbit question is resolved;
and every machine we do decide is, by the fact of admitting a certificate, not
one of the hard ones.  The principal results are impossibility theorems about
a certificate class, and their value is explanatory.  Each result in the paper
carries an explicit epistemic label---proved, proved by exhaustive finite
computation, machine-verified over a stated domain, or measured---and
Section~\ref{sec:ledger} tabulates all of them together with the threats to
validity.
\end{abstract}

\tableofcontents

\section{Introduction}

The $3x+1$ problem asks whether a two-branch affine map sends every positive
integer to $1$; it has been open for close to a century
\cite{lagarias-monthly,lagarias-overview}.  Its modern relatives are
\emph{Collatz-like halting problems}: small, fully explicit piecewise-affine
maps whose halting is equivalent to an orbit-avoidance question that resists
every known technique.  Conway showed that sufficiently general Collatz-like
maps simulate arbitrary computation \cite{conway72}, and Kurtz and Simon placed
the general halting problem for the class at $\Pi^0_2$-complete
\cite{kurtz-simon}.  Those reductions, however, require arbitrarily many
independently chosen affine branches; they say nothing about an individual
small map, nor about a parametrised family whose infinitely many branches are
generated by a fixed formula from finitely many integers.  This paper studies
one such family, chosen because the published halting-equivalent reduction of a
well-known open Turing machine is a member of it.

\subsection{Three kinds of object, and how they are related}
\label{sec:objects}

Confusion between the following is easy and consequential, so we fix
terminology first.

\paragraph{Turing machines.}
The busy-beaver community studies halting of concrete small Turing machines
\cite{bbchallenge-cryptids}.  The \textbf{Space Needle} is one: the $6$-state
machine
\[
\texttt{1RB1LA\_1LC0RE\_1LF1LD\_0RB0LA\_1RC1RE\_---0LD},
\]
found by mxdys in January 2025 \cite{bbchallenge-needle}.  Whether it halts is
open.

\paragraph{Collatz-like functions.}
Certain Turing machines admit a \emph{halting-equivalent reduction} to the
orbit of an arithmetic map.  For the Space Needle, Doucette's one-variable
reduction \cite{bbchallenge-needle} states: with $v(b)$ the $2$-adic valuation
of $b$ and $m$ its odd part, the machine halts if and only if the orbit of
$b_0=6$ under
\begin{equation}\label{eq:doucette}
b \;\longmapsto\; b + v(b) + \tfrac{3}{2}\bigl(m-1\bigr)
\qquad\text{(halting when $b$ is an exact power of two)}
\end{equation}
reaches an exact power of two.  A short computation identifies
\eqref{eq:doucette} as the member $(1,3,1,1,0)$ of
Definition~\ref{def:family}: writing $b=2^{v}(2k+1)$ gives
$b+v+3k=(2^{v+1}+3)k+(2^{v}+v)$.

\paragraph{Counter and rewrite systems.}
A third genre---machines given by a handful of rewrite rules on tuples of
counters---sits between the two.  The second machine studied here, which we
call \textbf{M3}, is of this kind: six rules on states $A(a,b)$, given in
Section~\ref{sec:baseq}, whose halting reduces to a base-$3$ orbit question.

\medskip
\noindent
\textbf{The consequence for reading this paper.}  Our census enumerates
\emph{Collatz-like functions}, not Turing machines: its members are five-tuples
of integers and were never presented as machines in the busy-beaver sense.
When we say a census member is ``decided'', we mean a proof that the orbit of
$x_0=3$ under that arithmetic map never reaches a power of two.  No Turing
machine's halting problem is decided anywhere in this paper.  The two
directions in which our results do touch genuine machines are both negative:
Corollary~\ref{cor:needle} and Theorem~\ref{thm:m3n1} say that a certain
proof technique can never settle the Space Needle's reduction, or M3.

\subsection{The research goals}
\label{sec:goals}

This paper reports part of a longer program.  The program's stated goals, in
the order in which they compound, are:

\begin{description}
\item[\textbf{G1} (observe).] Establish new structural facts about
  Collatz-like functions: normal forms, invariants, closed-form criteria.
\item[\textbf{G2} (decide).] Decide halting for specific machines wherever
  instance structure permits.
\item[\textbf{G3} (explain).] Sharpen the account of \emph{what makes the hard
  instances hard}---ideally by proving that named techniques cannot work,
  rather than by reporting that they have not worked.
\item[\textbf{G4} (transfer).] Grow a toolkit that moves between machines
  rather than being rebuilt per instance.
\end{description}

\subsection{Progress against the goals: summary}
\label{sec:summary}

Section~\ref{sec:goals-revisited} gives the assessment in full, with what did
not move.  In brief:

\begin{center}
\small
\begin{tabular}{@{}llp{8.1cm}@{}}
\toprule
goal & verdict & what this paper contributes \\
\midrule
G1 & substantial &
  the sieve lemma (Lemma~\ref{lem:universal}) and its linear corollary; the
  completeness theorem (Thm~\ref{thm:complete}); closed-form separation
  criteria at $m=3,4,6,8$ (Thm~\ref{thm:criteria}); the sieve-group rank
  dichotomy (Remark~\ref{rem:rank}), now tested on a machine we did not
  construct (Remark~\ref{rem:sheep-dichotomy}) \\
G2 & \textbf{moved, but not on the target class} &
  $318$ census members decided by certificate, $10$ more by the sieve---all
  of them Collatz-like \emph{functions}, none of them Turing machines, and
  none of them hard: a member with a certificate is by definition not a hard
  instance.  Neither flagship orbit question is resolved.  What is new here is
  a \emph{theorem about a wild cryptid} rather than a decision:
  Theorem~\ref{thm:sheep-sieve} and Corollary~\ref{cor:sheep-H} settle the
  arithmetic of the sheep machine completely, leaving only its orbit question
  \\
G3 & the paper's main yield &
  Theorem~\ref{thm:t15} and Corollary~\ref{cor:needle} replace a bounded
  search (``no certificate at $m\le 20{,}000$'') with an unconditional
  impossibility, \emph{and} identify the mechanism ($\delta$ sweeps every
  modulus; $\beta$'s parity governs the one $2$-adic escape hatch); plus the
  barriers of Section~\ref{sec:barriers} \\
G4 & demonstrated twice &
  the entire chain transfers to base $3$ (Section~\ref{sec:baseq}) and to a
  second \emph{wild} cryptid, the sheep machine (Section~\ref{sec:sheep});
  Lemma~\ref{lem:t16} shows the saturation argument needs neither the family
  nor a base; the toolkit's boundary is also mapped---two further machine
  families provably fall outside it \\
\bottomrule
\end{tabular}
\end{center}

\subsection{What this paper does not claim}
\label{sec:notclaimed}

We list these prominently because several are easy to over-read.

\begin{enumerate}
\item \textbf{No Turing machine is decided.}  The $328$ decided objects are
  arithmetic maps from a manufactured parameter box.  Nothing here reduces
  the number of open busy-beaver holdouts.
\item \textbf{Neither flagship is resolved.}  Whether the Space Needle's
  reduction, or M3, halts remains open.
\item \textbf{Decided $\Rightarrow$ not hard.}  A member admitting a
  separating congruence is, by that fact, outside the hard class.  Our G2
  contribution is a \emph{partition} of a family, not progress on its hard
  members.  Theorem~\ref{thm:t15} makes this sharper and less comfortable: the
  technique that decides $29.5\%$ of the census is provably unavailable for
  the flagships, forever.
\item \textbf{The impossibility theorems are about one certificate class.}
  They say nothing about non-halting proofs of other kinds; the natural next
  class (automatic invariants) is discussed in Section~\ref{sec:barriers} and
  is open.
\item \textbf{Completeness is relative to a box.}  Theorem~\ref{thm:complete}
  is a statement about the enumerated parameter box, not about all
  odd-$\beta$ members of the family.
\item \textbf{Universality of the family is open}, and
  Remark~\ref{rem:counting} shows one tempting route to settling it cannot
  work.
\end{enumerate}

\subsection{Epistemic conventions}
\label{sec:conventions}

Every numbered statement carries one of the following labels; the ledger in
Section~\ref{sec:ledger} lists all of them in one table.

\begin{description}
\item[\Proved] A mathematical proof is given here or is routine and
  indicated.  Where the proof contains a finite case analysis performed by
  computer, this is stated in the proof.
\item[\CAP] The statement is equivalent to a finite computation, and a
  computer performed it exhaustively.  This is a proof modulo correctness of
  the implementation; it is not a hand proof and we do not present it as one.
\item[\MV] Checked over an explicitly stated finite domain.  \emph{The general
  statement is not established.}  Such a claim is evidence, not a theorem.
\item[\Meas] An empirical quantity (a count, a frequency, a growth rate).
\item[\Interp] A judgment, comparison or placement---for instance a claim that
  our equations lie outside a known decidable fragment.  Not a formal result.
\end{description}

Each result is also accompanied by an \textbf{Evidence} note stating the exact
domain over which computations were run.  Threats to the validity of the
computational component are discussed in Section~\ref{sec:threats} and are not
negligible.

\section{The family}
\label{sec:family}

\begin{definition}[one-schema valuation machines]
\label{def:family}
A machine is a tuple $(\alpha,\beta,\gamma,\delta,\varepsilon)\in\Z^5$.  For
$x\in\N$ write $x=2^{v}u$ with $u$ odd, $v=\vp{2}(x)$.  If $u=1$ the machine
\textbf{halts}.  Otherwise $u=2k+1$ with $k\ge 1$, so $x=2^{v+1}k+2^{v}$, and
\[
F(x)=A_v\,k+B_v,\qquad
A_v=\alpha 2^{v+1}+\beta,\qquad
B_v=\gamma 2^{v}+\delta v+\varepsilon .
\]
The machine is \emph{well defined} if $F$ maps non-halting positive integers to
positive integers.
\end{definition}

The map reads the $2$-adic valuation of its argument, chooses affine
coefficients accordingly, and iterates; halting is landing exactly on a power
of two.  As a function of $x$ the branch has slope $A_v/2^{v+1}\to\alpha$, so
$\alpha$ is the asymptotic multiplier.

\begin{proposition}[identification of the Space Needle's reduction]
\label{prop:needle-id}
\Proved{}\; The member $(1,3,1,1,0)$ of Definition~\ref{def:family} coincides
with Doucette's reduction~\eqref{eq:doucette}.
\end{proposition}

\begin{proof}
With $b=2^{v}(2k+1)$ we have $m-1=2k$, so
$b+v+\frac32(m-1)=2^{v+1}k+2^{v}+v+3k=(2^{v+1}+3)k+(2^{v}+v)$, which is
$A_vk+B_v$ for $(\alpha,\beta,\gamma,\delta,\varepsilon)=(1,3,1,1,0)$.  The
halting conditions agree by definition.
\end{proof}

\begin{evidence}
Independently checked against an implementation of \eqref{eq:doucette}: zero
mismatches for $2\le x<300{,}000$; the orbit from $x_0=3$ is
$3,6,10,17,41,101,251,\dots$, which joins the published sequence at $6$.
\end{evidence}

\section{The census and exact congruence decision}
\label{sec:census}

\subsection{The decision procedure}

Fix a machine and a modulus $M\ge2$.  On branch $v$ both the source
$x=2^{v+1}k+2^{v}$ and the target $F(x)=A_vk+B_v$ are affine in $k$, so as $k$
ranges over $\Z_M$ the pair traces the graph of one affine correspondence
$\varphi_v$ on $\Z_M$.  Both $A_v$ and $B_v$ mod $M$ depend on $v$ only
through $(2^{v}\bmod M,\;v\bmod M)$, a state space of size at most $M^2$; so
iterating $v$ until that state repeats enumerates \emph{every} branch, and
$R_M=\bigcup_v\operatorname{graph}(\varphi_v)$ is exactly computable.  The
orbit of $x_0$ lies in the $R_M$-closure of $x_0\bmod M$, and the halting
values have residues forming a computable eventually-periodic set $H_M$.

\begin{definition}
$M$ \emph{separates} a machine from $x_0$ if the $R_M$-closure of $x_0\bmod M$
is disjoint from $H_M$.  Such an $M$ is a \emph{congruence certificate} and
proves non-halting from $x_0$.
\end{definition}

\begin{proposition}\label{prop:exact}
\Proved{}\; The procedure decides, for each $M$, whether $M$ separates; hence
it finds a separating modulus $\le M_{\max}$ iff one exists.  It is complete
for its class: unlike a test on a finite orbit prefix, the closure is exact.
\end{proposition}

\begin{proof}
Finiteness of the branch enumeration is the state-space argument above.  The
closure is taken over all $k$ including $k=0$, which only adds edges and hence
only enlarges the closure, so a disjointness verdict remains sound.
\end{proof}

\subsection{Census results}

\begin{theorem}[census]\label{thm:census}
\CAP{}\; In the box $\alpha\in\{1,2,3\}$, $\beta\in\{-1,\dots,7\}$,
$\gamma\in\{1,2,3\}$, $\delta\in\{0,1,2\}$, $\varepsilon\in\{-2,\dots,2\}$:
of $1{,}080$ tuples, $3$ are not well defined and $1{,}077$ are analysed.
Exactly $318$ ($29.5\%$) admit a separating congruence at some $M\le 64$ and
hence provably never halt from $x_0=3$.
\end{theorem}

\begin{evidence}
Every one of the $318$ certificates was re-audited by an independent
brute-force check over the integers ($x<200{,}000$, powers of two enumerated
completely, not reusing the branch enumeration that produced the certificate):
\textbf{zero failures}.  Smallest deciding moduli:
$3\,(41)$, $4\,(105)$, $6\,(68)$, $7\,(3)$, $8\,(18)$, $10\,(34)$, $12\,(17)$,
$14\,(2)$, $16\,(14)$, $20\,(6)$, $22\,(1)$, $24\,(2)$, $32\,(1)$, $36\,(1)$,
$40\,(2)$, $56\,(3)$.  \Meas{}: decidability concentrates at even $\beta$
($64$--$79\%$ decided) against odd $\beta$ ($3.7$--$8.1\%$);
Section~\ref{sec:saturation} proves the mechanism behind this.
\end{evidence}

\begin{theorem}[separation criteria at small moduli]\label{thm:criteria}
\CAP{}\; Separation at $M$ depends on the parameters only through their
residues mod $M$, so each criterion is a finite object on $(\Z_M)^5$.  Write
$A_0=2\alpha+\beta$, $B_0=\gamma+\varepsilon$, $T=A_0+B_0$.
\begin{enumerate}
\item[(T3)] Separation at $3$ $\iff$
  $\delta\equiv0,\ \gamma\equiv\alpha,\ \beta+\varepsilon\equiv0 \pmod 3$.
\item[(C4)] Separation at $4$ $\iff$ $\beta$ even and
  $2\alpha+\beta+\gamma+\varepsilon\equiv3\pmod4$.
\item[(C6)] Separation at $6$ $\iff$ one of: (a) $A_0$ even and
  $T\equiv3\pmod6$; (b) $A_0\equiv0,4\pmod6$ and $T\equiv5\pmod6$;
  (c) (T3) holds; (d) $A_0\equiv B_0\equiv0\pmod6$, $\delta\equiv0\pmod3$,
  $\gamma+\delta\equiv\beta+5\pmod6$.
\item[(C8)] At $8$ only branches $v\in\{0,1\}$ occur and $A_0$ must be even;
  separation is then determined by the forward orbit of $3$ under an explicit
  four-node map.
\end{enumerate}
\end{theorem}

\begin{proof}[Proof of (T3) and (C4), and status of (C6), (C8)]
Mod $3$ the powers of two are $\{1,2\}$, so a separating class is $\subseteq
\{0\}$ and contains $x_0=3$, hence equals $\{0\}$; then $3\mid x$ forces
$k\equiv1$ and
$F(x)\equiv(-1)^{v}(\gamma-\alpha)+\beta+\delta v+\varepsilon$, whose vanishing
at $v=0,1,2$ gives the three congruences and conversely.  Mod $4$ the powers of
two occupy $\{1,2,0\}$, so the only admissible class is $\{3\}\ni x_0$; only
$v=0$ has odd sources and $x\equiv3\pmod4$ iff $k$ is odd, giving targets
$A_0+B_0$ and $3A_0+B_0$, both $\equiv3$ iff $2A_0\equiv0$ and $T\equiv3$.
For (C6) and (C8) we give the case decomposition in the repository; the
statements as displayed are established by exhaustive verification on
$(\Z_M)^5$, which is a complete finite check, and we therefore label them
\CAP{} rather than \Proved{}.
\end{proof}

\begin{evidence}
Exhaustive on $(\Z_M)^5$: $9/243$ tuples separate at $3$; $128/1{,}024$ at $4$;
$1{,}308/7{,}776$ at $6$; $5{,}600/32{,}768$ at $8$; zero mismatches against
the decision procedure in every case.  Against the census: the criteria predict
the decided sets exactly ($41$, $105$, $68$, $18$ machines; zero discrepancies
in either direction).
\end{evidence}

\begin{remark}
The Space Needle's reduction has $\delta=1$ and so fails (T3) immediately.
Section~\ref{sec:saturation} upgrades this one-line observation into a theorem
covering every modulus.
\end{remark}

\section{The last-step sieve, ten machines, and completeness}
\label{sec:sieve}

Since halting is landing on a power of two, one may ask which branches can be
the \emph{last} step.  When branch $v$ is expanding, its affine fixed point
$P_0/Q_0$ gives the classical constraint: a halt can follow a $v$-step only if
$Q_0 2^{e}\equiv P_0 \pmod{A_v}$ is solvable, where
$P_0=(2B_v-A_v)2^{v}$ and $Q_0=2^{v+1}-A_v$.

\begin{lemma}[sieve lemma]\label{lem:universal}
\Proved{}\; Let $\beta$ be odd, so $A_v$ is odd.  For every member of the
family and every expanding branch $v$,
\[
\text{branch $v$ can immediately precede a halt}
\iff B_v\in\langle2\rangle \pmod{A_v},
\]
where $\langle2\rangle$ is the multiplicative group generated by $2$.
\end{lemma}

\begin{proof}
Reduce mod $A_v$: the $A_v$-terms vanish, leaving $Q_0\equiv2^{v+1}$, a unit,
and $P_0\equiv 2B_v2^{v}=B_v2^{v+1}$.  Hence $P_0Q_0^{-1}\equiv B_v$, and
solvability of $2^{e}\equiv P_0Q_0^{-1}$ is exactly the stated membership.
\end{proof}

\begin{evidence}
All $672$ well-defined odd-$\beta$ members of the box, every branch $v\le80$:
$51{,}030$ branches, zero mismatches against an independent implementation via
the fixed point.
\end{evidence}

\begin{corollary}\label{cor:linear}
\Proved{}\; Since $\alpha2^{v+1}\equiv-\beta\pmod{A_v}$,
$\alpha\cdot 2B_v\equiv 2\alpha\delta v+2\alpha\varepsilon-\beta\gamma
\pmod{A_v}$: the sieve target is affine in $v$.  The exponential appears only
in the modulus, never in the target.
\end{corollary}

Tractability of the membership question is governed by the size of
$\langle2\rangle \bmod A_v$.  Four $(\alpha,\beta)$ classes are
\emph{listable}---$\ord_{A_v}(2)$ grows linearly in $v$, so the powers of two
form a short explicit list---and the remainder are \emph{thin}, with $\ord$
growing exponentially on average and no elementary handle. \Meas{}

\begin{theorem}[ten machines]\label{thm:ten}
\Proved{}\; The following ten members never halt after their first step;
equivalently their halting starts are exactly the powers of two, which halt at
step $0$.  (Labels follow the program's numbering, in which T3 is the criterion
of Theorem~\ref{thm:criteria} and not a machine.)
\begin{center}
\small
\begin{tabular}{@{}llll@{}}
\toprule
machine & $2B_v \bmod A_v$ & why every branch is forbidden & label\\
\midrule
$(1,1,2,0,1)$ & $0$ & $B_v=A_v$; $0$ is not a unit & T1\\
$(2,-1,2,0,1)$ & $3$ & odd, $1<3<A_v$ for $v\ge1$ & T9\\
$(2,-1,2,1,1)$ & $2v+3$ & odd, in range for $v\ge1$ & T2\\
$(2,-1,2,2,1)$ & $4v+3$ & odd, in range for $v\ge2$ & T4\\
$(2,-1,3,0,0)$ & $2^{v+1}+1$ & odd, in range for $v\ge1$ & T5\\
$(2,-1,3,1,0)$ & $2^{v+1}+2v+1$ & odd, in range for $v\ge2$ & T6\\
$(2,-1,1,2,-1)$ & $2^{v+1}+4v-2$ & $=2\cdot\text{odd}<A_v$, so
  $2^{v}=2-2v$: impossible & T10\\
$(2,1,1,2,-1)$ & $2^{v+1}+4v-2$ & signed list; $4v-2=2^{v+1}$: impossible
  & T7\\
$(3,3,2,0,1)$ & $\equiv0 \bmod d_v$ & $B_v=d_v=2^{v+1}+1\mid A_v$; powers of
  $2$ are units mod $d_v$ & T8\\
$(3,3,3,0,0)$ & $-3$ & $3\mid A_v$, $2B_v=A_v-3$; powers of $2$ mod $3$ are
  $\{1,2\}$ & T11\\
\bottomrule
\end{tabular}
\end{center}
In each row the second column is an exact integer identity (e.g.\
$2B_v=A_v+3$ for T9), the powers of two mod $A_v$ form the explicit listable
set of that machine's class, and the stated parity, size or divisor argument
excludes membership for all $v$ beyond finitely many boundary branches, which
land on $0$.
\end{theorem}

\begin{evidence}
Closed forms asserted as integer identities to $v\le4{,}000$ on a second,
independently written implementation; every branch $v\le200$ confirmed
forbidden; brute force on all ten (every non-power-of-two start below
$20{,}000$, $400$ steps): zero halts.  Two of the ten upgrade earlier
certificates from ``never halts from $x_0=3$'' to ``from any start''; T10 has a
genuine cycle ($F(3)=3$) and still never halts from any non-power-of-two start.
\end{evidence}

\begin{theorem}[completeness of the sieve method within the box]
\label{thm:complete}
\CAP{}\; Within the enumerated box, the ten machines of
Theorem~\ref{thm:ten} are the only ones all of whose branches are forbidden.
For every other well-defined odd-$\beta$ member there is an explicit
\emph{surviving-branch certificate}: a branch $v\le22$ and an exponent $e$
with $2^{e}\equiv 2B_v\pmod{A_v}$.
\end{theorem}

\begin{evidence}
Of $672$ machines: $10$ all-forbidden, $620$ with a surviving-branch
certificate (each verified by modular exponentiation), $42$ sieve-silent (the
$(1,-1)$ class, which has no expanding branch).
\end{evidence}

A surviving branch is a positive witness, not a failure to find one; the
method is therefore closed by a finite proof rather than by an exhausted
budget.  We stress two limits: the statement is relative to the box, and a
surviving branch does not make a machine halt---it certifies only that this
method cannot decide it.

For the Space Needle's reduction the lemma specialises to a clean open
condition.  With $A_v=2^{v+1}+3$ we get $2^{v+1}\equiv-3$, so
$\langle2\rangle=\{\pm3^{a}2^{i}\}$, and branch $v$ can precede a halt iff
\begin{equation}\label{eq:needle}
2v-3 \;\in\; \{\pm 3^{a}2^{i}\} \pmod{2^{v+1}+3}.
\end{equation}

\begin{observation}\label{obs:nothinning}
\Meas{}\; The set of branches satisfying \eqref{eq:needle} does not appear to
thin out: $18$ of the branches $v\le40$ survive, including $v=34,37,38$.  The
weighted forbidden mass is $0.2868$ and $19$ of the first $35$ valuations are
forbidden, independently reproducing an earlier measurement.  \emph{We draw no
asymptotic conclusion}; this is stated to record that no all-branches theorem
is available for the Needle by this route.
\end{observation}

\section{No congruence certificate: the saturation theorems}
\label{sec:saturation}

We now show the certificate class of Section~\ref{sec:census} is empty for the
flagships---not below a bound, but outright.  Write $M=2^{s}M'$ with $M'$ odd.
For odd modulus the branch correspondence is a single affine map
$\varphi_v(c)=\mu_vc+a_v$, and a short computation gives
$\mu_v=\alpha+\beta 2^{-(v+1)}$ and $a_v=B_v-A_v/2=(\gamma-\alpha)2^{v}+\delta
v+(\varepsilon-\beta/2)$.

\begin{theorem}[$\delta$-saturation]\label{thm:t12}
\Proved{}\; Let $M$ be odd, $\gcd(\delta,M)=1$ and $\gcd(\ord_M(2),M)=1$.  Then
for every residue $c$ the one-step image $\{\varphi_v(c):v\ge0\}$ is all of
$\Z_M$.
\end{theorem}

\begin{proof}
On a class $v\equiv v_0 \pmod{\ord_M(2)}$ the quantity $2^{v}$ is constant, so
$\mu_v$ and $(\gamma-\alpha)2^{v}$ are constant and the only surviving
$v$-dependence is $\delta v$: thus $a_v-a_{v_0}\equiv\delta(v-v_0)$.  As $v$
runs over the infinite progression $v_0+\ord_M(2)\Z$, the difference sweeps the
subgroup generated by $\gcd(\ord_M(2),M)=1$, i.e.\ all of $\Z_M$, and
$\gcd(\delta,M)=1$ preserves this.
\end{proof}

Since $\ord_p(2)\mid p-1$, the order hypothesis holds automatically at every
odd prime; already this shows no odd prime separates the Needle's reduction.
The hypothesis is scaffolding and is removed next.

\begin{lemma}[descent]\label{lem:descent}
\Proved{}\; Let $M$ be odd with $\gcd(\delta,M)=1$.  Then the conclusion of
Theorem~\ref{thm:t12} holds with no hypothesis on $\ord_M(2)$.
\end{lemma}

\begin{proof}
Set $g_0=M$, $g_{j+1}=\gcd(\ord_{g_j}(2),g_j)$, and let $I_j$ be the image of
$\{\varphi_v(c)\}$ in $\Z_{g_j}$.  For odd $g\ge3$, $\ord_g(2)\le\lambda(g)<g$,
so the chain strictly decreases to $1$.  Mod $g_j$ with $h=\ord_{g_j}(2)$: on
each class $v\equiv w\pmod h$ everything but $\delta v$ is constant, and
$\delta v$ sweeps the \emph{full} subgroup $\langle g_{j+1}\rangle$---here the
infinitude of the progression is essential, one period being insufficient.  So
$I_j$ is a union of $h$ cosets of $\langle g_{j+1}\rangle$, and reducing mod
$g_{j+1}$ annihilates the sweep and leaves exactly the base points, whence
$I_j=\Z_{g_j}\iff I_{j+1}=\Z_{g_{j+1}}$.  The chain bottoms out at $g=1$, where
fullness is trivial; walking back up gives $I_0=\Z_M$.
\end{proof}

\begin{evidence}
$2{,}500$ (machine, odd $M$, $c$) instances with $\gcd(\delta,M)=1$, of which
$708$ have $\gcd(\ord_M(2),M)>1$---exactly the cases the scaffold hypothesis
excluded---the one-step image fell short of $\Z_M$ in none.  The chain
strictly decreases for every odd $g<6{,}000$, longest chain length $7$.
\end{evidence}

\subsection{Even moduli: the parity escape hatch}
\label{sec:even}

A residue $c\bmod 2^{s}$ with $\vp{2}(c)=v<s$ \emph{pins the branch index}: the
low bits of the value are the valuation.  That is precisely the information an
odd modulus cannot carry, and it is the entirety of what an even modulus buys.
On a source $\equiv0\pmod{2^{s}}$ every branch $v\ge s$ is available, $k$ is
free mod $2^{s}$, and since $v+1>s$,
\[
F(x)\equiv \beta k+\delta v+\varepsilon \pmod{2^{s}} .
\]
If $\beta$ is odd, $\beta k$ sweeps all of $\Z_{2^{s}}$ and the $2$-adic
information is destroyed in one step; if $\beta$ is even only
$\langle\gcd(\beta,2^{s})\rangle$ is reached---the escape hatch.

\begin{theorem}[even-modulus saturation]\label{thm:t14}
\Proved{}\; Let $M=2^{s}M'$, $\beta$ odd, $\gcd(\delta,M')=1$, and suppose the
$R_M$-closure of $x_0$ contains some $c_0\equiv0\pmod{2^{s}}$.  Then that
closure is all of $\Z_M$.
\end{theorem}

\begin{proof}
From $c_0$ every branch $v\ge s$ is available (choose an odd representative mod
$M'$, possible since $M'$ is odd and $2^{v}$ is invertible there).  Splitting
the source congruence by CRT: mod $2^{s}$ it holds identically, so $k$ is
unconstrained; mod $M'$ it determines $k$.  The target's $2$-part sweeps
$\Z_{2^{s}}$ by the display above; its $M'$-part, over $v\ge s$, is the
$\delta$-sweep of Lemma~\ref{lem:descent} restricted to $v\ge s$, which costs
nothing since each order class contains infinitely many such $v$.  The product
is $\Z_{2^{s}}\times\Z_{M'}=\Z_M$.
\end{proof}

\begin{evidence}
Proof steps checked on $3{,}000$ instances; the conclusion on $7{,}840$
(odd-$\beta$ machine, even modulus) pairs: zero failures.  \emph{Falsifier}:
each of the census's $274$ even-modulus certificates must violate a
hypothesis---$267$ have $\beta$ even (they spend exactly the hatch) and $7$
have $\delta=0$; none would refute.  The $44$ odd-modulus certificates all
violate $\gcd(\delta,M)=1$.
\end{evidence}

\subsection{The reach lemma}

It remains to supply the $c_0$ of Theorem~\ref{thm:t14}.

\begin{lemma}[reach]\label{lem:reach}
\Proved{}\; Let $\beta$ be odd and $A_v>0$ for all $v$.  For every residue $c$
mod $M=2^{s}M'$ there is a lift $x\equiv c\pmod M$ whose $F$-orbit reaches a
value $\equiv0\pmod{2^{s}}$ within $s+1$ steps.
\end{lemma}

\begin{proof}[Proof (exponent ledger)]
If $c\equiv0$ there is nothing to prove; else let
$v_0=\vp{2}(c\bmod 2^{s})<s$ and take $x=c+2^{s}M'T$ with $T$ free.  Iterating
$F$ \emph{exactly} keeps the $T$-dependence affine, $y_n=p_n+q_nT$, with
$q_n=A_{v_{n-1}}\cdots A_{v_0}M'2^{e_n}$ and the ledger $e_1=s-(v_0+1)$,
$e_{n+1}=e_n-(v_n+1)$: forming $k=(y_n-2^{v_n})/2^{v_n+1}$ divides the
$T$-coefficient by $2^{v_n+1}$, and multiplying by the odd $A_{v_n}$ changes
nothing $2$-adically.  Two events are exhaustive.  \emph{Done-event}: some
$p_n$ is divisible by $2^{e_n}$, and then $y_n\equiv0\pmod{2^{s}}$ has a
solution $T\equiv-(p_n/2^{e_n})(q_n/2^{e_n})^{-1}\pmod{2^{s-e_n}}$, which may
be taken arbitrarily large.  \emph{Ledger-event}: $e$ reaches $0$, so $k$ is a
unit times $T$ mod $2^{s}$ and, $A_v$ being odd, $y_{n+1}=A_vk+B_v$ takes every
residue mod $2^{s}$ including $0$.  While neither has fired,
$\vp{2}(p_n)<e_n$ makes $v_n$ independent of $T$, so the branch sequence is
well defined, and the ledger strictly decreases from at most $s-1$; the process
ends within $s+1$ steps.  For $T$ large every $y_n$ is positive with odd part
exceeding $1$, so no step halts en route.
\end{proof}

\begin{evidence}
Constructive: the recipe was used to \emph{build} the witness lift and its
orbit was \emph{run}, on $4{,}000$ (machine, modulus, residue) instances---a
hit on $0\bmod 2^{s}$ every time, deepest at $9$ steps against the bound $10$.
The ledger bookkeeping was checked in isolation on $3{,}000$ chains and the
$T$-independence of the branch sequence confirmed.  \emph{Falsifier}:
$\beta$-even machines with even-modulus certificates are $2$-adically confined,
so reach must fail for them---it does, in $40$ of $40$ tested.
\end{evidence}

\begin{theorem}[no congruence certificate]\label{thm:t15}
\Proved{}\; Let a member have $\beta$ odd, $A_v>0$, and $\gcd(\delta,M')=1$
where $M'$ is the odd part of $M$.  Then the $R_M$-closure of every residue is
all of $\Z_M$: no modulus $M$ separates, and no branch can be excluded.
\end{theorem}

\begin{proof}
Odd $M$ is Lemma~\ref{lem:descent}.  For even $M$, Lemma~\ref{lem:reach} places
a residue $\equiv0\pmod{2^{s}}$ in the closure and Theorem~\ref{thm:t14}, with
Lemma~\ref{lem:descent} supplying its $M'$-part, finishes.
\end{proof}

\begin{corollary}\label{cor:needle}
\Proved{}\; \textbf{No congruence certificate at any modulus separates the
Space Needle's reduction.}  ($\beta=3$ odd, $\delta=1$, $A_v=2^{v+1}+3>0$, so
every hypothesis holds at every $M$.)  Equivalently: the technique that decides
$29.5\%$ of the census can never decide this instance.
\end{corollary}

\begin{evidence}
$7{,}616$ (odd-$\beta$ machine, modulus) pairs meeting the hypotheses: closure
short of $\Z_M$ in none.  The Needle's reduction directly at every modulus
$2\le M\le400$: saturated in all cases.
\end{evidence}

What began as ``$\delta$ blocks the cheapest certificate'' (the one-line
failure of (T3)) is now: $\delta$ blocks \emph{every} certificate at
\emph{every} modulus, and the escape hatch that accounts for $267$ of the
$274$ even-modulus certificates is closed exactly by $\beta$ odd.

\section{Transfer to base $q$: machine M3}
\label{sec:baseq}

None of the proofs above uses base $2$ essentially.  We demonstrate the
transfer on a second machine.  \textbf{M3} is the six-rule system
\[
\begin{aligned}
&A(1,3k)\to\text{HALT},\quad A(1,3k{+}1)\to A(3k{+}4,1),\quad
A(1,3k{+}2)\to A(3k{+}3,2),\\
&A(3k,b)\to A(k,b{+}2k{+}1),\quad A(3k{+}1,b)\to A(4k{+}b{+}3,1),\quad
A(3k{+}2,b)\to A(4k{+}b{+}5,1),
\end{aligned}
\]
started at $A(1,1)$; its halting is equivalent to the orbit of the return map
to the $b=1$ section reaching a power of $27$, which is open.

\begin{proposition}[M3's branch table]\label{prop:m3table}
\MV{}\; Writing a non-spine value as $a=3^{j+1}m+r3^{j}$ with $j=\vp{3}(a)$ and
leading digit $r\in\{1,2\}$ (the spine being the pure powers $3^{j}$, which
halt iff $3\mid j$), the return map is
\[
G(a)=A_j m+B_{j,r},\qquad A_j=3^{j+1}+1,\qquad
B_{j,r}=r3^{j}+j+c_r,\ \ c_1=3,\ c_2=4,
\]
reached in $j+1$ raw steps.  In the coordinates of
Definition~\ref{def:family} read base $3$: $\alpha=1$, $\beta=1$,
$\gamma=r$, $\delta=1$, $\varepsilon_r=r+2$.
\end{proposition}

\begin{evidence}
Checked against the raw six rules on $59{,}988$ non-spine values $a<60{,}000$:
zero mismatches, including the step count $j+1$; spine behaviour exact for
$j=1,\dots,11$.  \emph{We label this \MV{} deliberately}: the induction over
the $j$ divide-steps is routine but is not written out here, and
Theorems~\ref{thm:m3n2} and~\ref{thm:m3n1} below inherit that dependency.
\end{evidence}

\begin{lemma}[base-$q$ sieve lemma]\label{lem:baseq}
\Proved{}\; For a branch sending $x=q^{w}m+s$ to $Am+B$ with $\gcd(q,A)=1$ and
halting set $H=\{c\rho^{t}\}$, $\gcd(\rho,A)=1$: the branch can immediately
precede a halt iff $B\in c\langle\rho\rangle\pmod A$.  For M3 the fixed point
gives $Q_0=3^{j+1}-A_j=-1$ exactly, so the condition reads
$27^{t}\equiv B_{j,r}\pmod{A_j}$.
\end{lemma}

The base-$2$ form of Lemma~\ref{lem:universal} (with its factor $2$) was a
normalisation artifact; the invariant content in every base is
\emph{target}~$\equiv B$.

\begin{theorem}[M3: last two steps pinned]\label{thm:m3n2}
\CAP{}\; For every branch index $j$: a halt of M3 can only follow a branch with
$\vp{3}=1$, and the same holds one step earlier.  The excluded-pair frequency
$1-(2/9)^2=77/81$ is exact.  Both surviving depth-one branches are realised:
$21\to27$ and $478293\to27^{4}$.
\end{theorem}

\begin{evidence}
Depth-one verdicts agree with an independent solver on $50$ branches, survivors
exactly $(1,1)$ and $(1,2)$.  The depth-two argument (a parity lemma removing
even $j$; for odd $j\ge5$ the targets lie strictly inside the gaps of
$\{\pm3^{i}\}$; small $j$ checked exactly) was machine-checked to $j=2{,}000$
and cross-checked against an exact finite test for $2\le j\le500$.
\emph{Caveat}: the present author's independent re-check of this theorem
covered only its brute-force consequence (among all non-spine $a<10^{6}$,
exactly three starts halt within two return steps---$15$, $21$,
$478293$---with every branch value en route having $\vp{3}=1$), not the
depth-two closed forms.  This is the least independently corroborated result in
the paper.
\end{evidence}

\begin{theorem}[M3: no congruence certificate]\label{thm:m3n1}
\Proved{} (modulo Proposition~\ref{prop:m3table})\; For every modulus $M\ge2$
the closure of every residue under M3's branch relation is all of $\Z_M$; no
congruence certificate at any modulus separates M3's orbit from the residues of
$\{27^{k}\}$.
\end{theorem}

\begin{proof}
The chain of Section~\ref{sec:saturation} with $2\to3$: ``odd $M$'' becomes
``$M$ coprime to $3$''; ``$\beta$ odd'' becomes ``$\beta$ coprime to $3$''
(here $\beta=1$); the exponent ledger spends trits.  M3 has $\delta=1$, so all
hypotheses hold at every modulus.  Spine transitions only enlarge closures.
\end{proof}

\begin{evidence}
Independent implementation: every modulus $2\le M\le250$, three start residues
each ($747$ closures), all equal to $\Z_M$.
\end{evidence}

\begin{remark}[the flagship dichotomy]\label{rem:rank}
\Proved{}\; Both flagships satisfy Theorem~\ref{thm:t15}'s hypotheses for the
same reason ($\delta=1$, $\beta$ coprime to the base) and so resist every
congruence certificate.  What separates their fates under the \emph{sieve} is
one class parameter, the rank of $\langle q\rangle\bmod A_v$: M3's class forces
$q^{j+1}\equiv-1$, making $\langle3\rangle=\{\pm3^{i}\}$ listable (rank one), so
the gap argument closes every branch question; the Needle's class gives
$\langle2\rangle=\{\pm3^{a}2^{i}\}$ (rank two, thin), which is the open
condition~\eqref{eq:needle}.  In the census, $120$ members share the Needle's
group.
\end{remark}

\begin{remark}[interface boundary]\label{rem:boundary}
\Interp{}\; The toolkit requires a valuation-indexed branch family together
with a halting set defined by a condition on the \emph{value}.  Two further
machine families studied in the program---the Hydra/Antihydra maps and the
FRACTRAN cryptid Fenrir---fall outside it: their branch modulus is constant and
halting is a condition on a path counter rather than on the value, so the sieve
has no object to apply to.  We record this as a mapped boundary, not a failure.
\end{remark}

\section{Two tests outside the design set}
\label{sec:wild}

Everything above was developed against the Space Needle's reduction.  That is
a real threat to validity, and this section reports the two tests we can run
against it: a second \emph{wild} cryptid that lies inside the interface, and
an old case-file machine that lies outside the family but inside the
saturation argument.

\subsection{The sheep machine: a second wild instance}
\label{sec:sheep}

The Turing machine \verb|1RB1LA_0LC0RC_1LE1RD_1RE1RC_1LF0LA_---1LE| was found
by \emph{sheep} on 7 April 2026 and is listed as a Cryptid on the bbchallenge
wiki, which states the halting-equivalent reduction
\[
f(n)=\begin{cases}
\text{HALT} & \oddpart(n)=1,\\
n+\vp{2}(n)+3 & \oddpart(n)=3,\\
n+\vp{2}(n)+\tfrac{\oddpart(n)-1}{2} & \oddpart(n)>3,
\end{cases}
\qquad\text{started at } n=5,
\]
together with two lemmas: that $m2^a+a+\frac{m-1}{2}$ is never a power of two
for $a\ge2$, $m\ge5$; and that $3\cdot2^a+a+3$ is never a power of two.

\begin{proposition}[identification]\label{prop:sheep-id}
\Proved\; On the generic branch ($\oddpart(n)>3$), writing $n=2^v m$ with
$m=2k+1$,
\[
f(n)=A_vk+B_v,\qquad A_v=2^{v+1}+1,\quad B_v=2^v+v,
\]
so the generic branch is the census member $(\alpha,\beta,\gamma,\delta,
\varepsilon)=(1,1,1,1,0)$.  The Needle's reduction is $(1,3,1,1,0)$: the two
machines differ in $\beta$ alone.
\end{proposition}

\emph{Evidence.}  $15{,}920$ systematic and $200{,}000$ random $(v,k)$ pairs
with $v<200$, $k<10^{30}$: no mismatch.  For $n<60{,}000$, $f$ differs from
the census member exactly on the values with $\oddpart(n)=3$.

The exceptional branch is not decoration.  The census records member
$(1,1,1,1,0)$ as \emph{halting}: $3\mapsto4=2^2$.  The sheep machine's third
branch intercepts precisely that value, sending $3\mapsto6$.  One extra branch
on one odd part converts a halting census member into an open problem---the
sharpest statement we can make about what a manufactured box misses.

\begin{theorem}[the sieve closes]\label{thm:sheep-sieve}
\Proved\; For the sheep machine, branch $v$ can immediately precede a halt if
and only if $v\in\{0,1\}$.
\end{theorem}

\begin{proof}
A generic step lands on a power of two iff $2^t\equiv B_v \pmod{A_v}$ for some
$t$, i.e.\ iff $B_v\in\langle2\rangle$.  Since $A_v=2^{v+1}+1$ we have
$2^{v+1}\equiv-1$ and $2^{-1}\equiv2^v+1$, so $B_v=2^{-1}-1+v$ and the
condition is
\[
2^{t+1}\equiv 2v-1 \pmod{2^{v+1}+1}.
\]
Also $2^{v+1}\equiv-1$ forces
$\langle2\rangle\subseteq\{2^i\}_{i\le v}\cup\{A_v-2^i\}_{i\le v}$.  If
$2v-1=2^i$ then $2v-1$ odd gives $i=0$, hence $v=1$.  If $2v-1=A_v-2^i$ then
$2^i=2^{v+1}-2v+2$, which for $v\ge3$ lies strictly between $2^v$ and
$2^{v+1}$ (as $2^v>2v-2$ and $2v>2$) and so is not a power of two; $v=2$ gives
$6$.  Finally $v=0$: $A_0=3$ and $2v-1\equiv2=2^1\in\langle2\rangle$.
\end{proof}

The wiki's first lemma is the $v\ge2$ half of Theorem~\ref{thm:sheep-sieve}
with the threshold $a\ge2$ assumed; here the threshold is a conclusion, and
the two surviving branches are located at the same time.  The second lemma is
recovered the same way: for $a\ge3$, $3\cdot2^a+a+3$ lies strictly between
$3\cdot2^a$ and $2^{a+2}$, and $a\le2$ gives $6,10,17$.

\begin{corollary}[the halting set in closed form]\label{cor:sheep-H}
\Proved\; A value halts in one step iff it is a power of two, and reaches a
power of two in exactly one step iff it lies in
\[
H_0=\Bigl\{\tfrac{2^{2j+1}+1}{3}\Bigr\}_{j\ge2}=11,43,171,683,\dots
\quad\text{or}\quad
H_1=\Bigl\{\tfrac{2(2^{4j}-1)}{5}\Bigr\}_{j\ge2}=102,1638,26214,\dots
\]
so $H=\{2^i\}\cup H_0\cup H_1$.
\end{corollary}

\emph{Evidence.}  Brute force over $n<300{,}000$ returns exactly the $11$
predicted members.  The excluded first member of each family is $m=3$ in both
cases---the exceptional branch again.

Finally $\beta=1$ is odd and $\delta=1$, so Theorem~\ref{thm:t15} applies
verbatim: no congruence certificate at any modulus separates the sheep
machine.  \emph{Evidence.}  Direct closure computation gives all of $\Z_M$
for every $M\le200$.

\begin{remark}[the dichotomy, tested]\label{rem:sheep-dichotomy}
\Interp\; Remark~\ref{rem:rank} predicts that the last-step sieve closes when
$\langle2\rangle$ modulo $A_v$ is a rank-one listable group and stalls when it
is rank two and thin.  The sheep machine has $A_v=2^{v+1}+1$, hence
$\langle2\rangle=\{\pm2^i\}$, rank one; the Needle has $A_v=2^{v+1}+3$, hence
$\{\pm3^a2^i\}$, rank two.  The sieve closes for the former and not the
latter.  This is the first test of that prediction on a machine we did not
construct, and it was on record before the sheep machine was examined.
\end{remark}

The result is a cryptid whose \emph{arithmetic} is settled completely---we
know exactly which values can halt, in closed form---and whose orbit question
is untouched.  The Needle is one where neither is settled.  We regard this as
the paper's clearest evidence that the sieve-group rank, and not the machine,
is what governs the last-step analysis.

\subsubsection*{Depth two, and why depth does not help}

The same reduction runs one step further.  Since $3N=2^{2j+1}+1$ and
$5N=2^{4j+1}-2$ identically, ``$N\equiv B_v \pmod{A_v}$'' clears denominators
into $2^{2j+1}\equiv 3B_v-1 \pmod{3A_v}$ and $2^{4j+1}\equiv 5B_v+2
\pmod{5A_v}$; the larger modulus is required because $3\mid A_v$ whenever $v$
is even.

\begin{theorem}[depth two]\label{thm:sheep-d2}
\Proved\; No branch $v\ge6$ can be the second-to-last step before a halt; the
depth-two survivor set is exactly $\{0,1,2,3,5\}$.
\end{theorem}

\begin{proof}
Reduce each condition modulo $A_v$ (a necessary condition).  Because
$2^{v+1}\equiv-1$, $\ord_{A_v}(2)$ is even, so the exponent parity of each
element of $\langle2\rangle=\{2^i\}_{i\le v}\cup\{A_v-2^i\}_{i\le v}$ is well
defined, and the two targets reduce to $2^v+3v-2$ and $2^v+5v$.  Each must be
$\pm2^i$.  For $H_0$: $2^v<2^v+3v-2<2^{v+1}$ kills the first case for
$v\ge4$, and the second gives $2^i=2^v-3v+3\in(2^{v-1},2^v)$, impossible for
$v\ge5$.  For $H_1$: $2^v<2^v+5v<2^{v+1}$ for $v\ge5$, and $2^i=2^v-5v+1$ lies
in $(2^{v-1},2^v)$ for $v\ge6$---while $v=5$ gives $32-25+1=8=2^3$, a genuine
solution.  Branches $v\le5$ are then checked directly.
\end{proof}

Iterating is possible in closed form.  Call $N(i)=(2^{\alpha+ei}+b)/c$ with
$c$ odd a \emph{geometric family}; the halting set is one such family, and the
shape is closed under preimages, since $f(x)=N(i)$ for $x=2^v(2k+1)$ reduces
to $2^{\alpha+ei}\equiv cB_v-b \pmod{cA_v}$, whose solutions form a class
$i\equiv i_0 \pmod P$, and substituting back yields a family with $e'=eP$,
$c'=cA_v$.  So each depth is computed exactly, with no search over $x$.

\begin{center}
\small
\begin{tabular}{@{}clrrr@{}}
\toprule
depth & survivor set & families & forbidden branch mass & admissible word mass\\
\midrule
1 & $\{0,1\}$ & 2 & $0.250000$ & $0.750000$\\
2 & $\{0,1,2,3,5\}$ & 7 & $0.046875$ & $0.714844$\\
3 & $\{0,\dots,6\}$ & 23 & $0.007812$ & $0.709259$\\
4 & $\{0,\dots,6,9\}$ & 90 & $0.006836$ & $0.704411$\\
5 & $\{0,\dots,10,13\}$ & 346 & $0.000427$ & $0.704110$\\
6 & $\{0,\dots,13\}$ & 1421 & $0.000061$ & $0.704067$\\
\bottomrule
\end{tabular}
\end{center}

\begin{remark}[the sieve saturates]\label{rem:sheep-ladder}
\Meas\; The forbidden mass per depth collapses geometrically while the family
count grows by a factor $\approx3.8$ per depth, so the admissible word mass is
a \emph{convergent} product, measured at $0.704067$ by depth six and
extrapolating to $\approx0.70406$.  The last-step sieve, at any depth, can
therefore never forbid more than about $29.6\%$ of branch words.  We note that
the depth-five data alone extrapolated to $0.704045$, and depth six then
measured $0.704067$; the extrapolation was accurate to five decimals before
the check existed.  Depth seven was not attempted: cost grows by a factor
$\approx15$ per depth ($1$\,s, $80$\,s, $6.4$\,h for depths four to six).
\end{remark}

The mechanism is legible: the depth-$d$ target is a union of $F_d$ families,
so branch $v$ survives if \emph{any} is reachable; each is reachable with
probability $\approx\ord(2)/A_v=O(v)/2^v$, so $v$ survives up to
$v\approx\log_2 F_d$, which grows linearly in $d$.  Survivor sets grow, the
forbidden mass decays geometrically, and the product converges.  \emph{Depth
is not a resource that buys arbitrarily much}---which, with
Theorem~\ref{thm:t15}, closes both congruence-flavoured lanes on this machine
with numbers attached, and leaves automatic invariants as the only one
untried.  Evidence: ground truth by brute force to $3\times10^6$ agrees at
depths one to four; the recursion is verified against $f$ itself on two
members of every family generated.

\subsection{A machine-independent saturation lemma}
\label{sec:t16}

The descent of Lemma~\ref{lem:descent} does not need the family at all.

\begin{lemma}[linear-exponential saturation]\label{lem:t16}
\Proved\; Let $M\ge1$ and let $A,\delta,\varepsilon,\kappa,\rho$ be integers
with $\gcd(\delta,M)=\gcd(\rho,M)=1$.  Put $f(n)=\delta n+\varepsilon+
\kappa\rho^n$ and $R=\{(c,\,Ac+f(n))\bmod M: n\ge n_0\}$.  Then from every
$c\in\Z_M$, the residues reachable in at least one $R$-step are all of $\Z_M$.
There is no hypothesis on $A$, none on $\kappa$, and $n_0$ is irrelevant.
\end{lemma}

\begin{proof}
Strong induction on $M$; $M=1$ is trivial.  Let $T=\ord_M(\rho)$ and
$g=\gcd(T,M)$; as $T\le\varphi(M)<M$ we have $g<M$.  For any $c$ and $n_1$,
the successors $\{Ac+f(n_1+jT)\}_{j\ge0}=Ac+f(n_1)+\delta T\Z=
Ac+f(n_1)+\langle g\rangle$, using $\gcd(\delta,M)=1$; so every residue
reached in $\ge1$ step lies in a full coset of $\langle g\rangle$ that is also
reached.  If $g=1$ we are done.  Otherwise $\pi:\Z_M\to\Z_M/\langle g\rangle
\cong\Z_g$ carries $R$ to a relation of the same shape with
$\gcd(\delta,g)=\gcd(\rho,g)=1$; by induction $\pi$ of the reachable set is
all of $\Z_g$, so the reachable set meets every coset and hence, by the
previous sentence, contains each of them.
\end{proof}

\emph{Evidence.}  $1{,}705$ random parameter tuples with $M<400$ all saturate;
dropping $\gcd(\delta,M)=1$ makes $437$ of $671$ instances fail, so the
hypothesis is load-bearing.

Applying it to a machine outside the family: the counter system \textbf{M1}
studied earlier in this project reduces to a return map $F$ whose dominant
branch word satisfies, for $n\ge1$ and exactly the $D$ with
$16\cdot2^n-2n-10\le D\le 20\cdot2^n-2n-13$,
\[
F(D)=16D-240\cdot2^n+32n+169 \tag{\Proved}
\]
(both interval endpoints derived: the lower from the exit guard, the upper
from $n$ being the cascade length).  That interval has length $4\cdot2^n-2$,
so once $2^n\ge(M+2)/4$ it contains a complete residue system modulo $M$, and
therefore the true edge relation of $F$ modulo $M$ contains the relation of
Lemma~\ref{lem:t16} with $A=16$, $\delta=32$, $\varepsilon=169$,
$\kappa=-240$, $\rho=2$.

\begin{theorem}[M1]\label{thm:m1n1}
\Proved\; No congruence certificate at any odd modulus proves M1 non-halting;
this already follows from its single dominant branch word.
\end{theorem}

\emph{Evidence.}  Verified exhaustively for every odd $M\le401$ and for $25$
random odd $M<1200$.  On the even side the orbit is confined to one class
modulo $16$ and no $2^e$ with $e\ge5$ says more, so the complete congruence
content of M1 is a single residue class---which does not separate it from its
halting set.  M1 is thus the first machine we place on the $\beta$-even side
of the parity dichotomy of Section~\ref{sec:even}.

\section{Expressiveness barriers}
\label{sec:barriers}

Whether the family is Turing-complete is open.  This section bounds what such
a proof must overcome and refutes one tempting genre of impossibility argument.

\begin{theorem}[branch rigidity]\label{thm:rigid}
\Proved{}\; Branch $v$ sends every non-halting input to the same valuation iff
$\vp{2}(B_v)<\vp{2}(A_v)$, and then $\vp{2}(F(x))=\vp{2}(B_v)$ throughout.  If
this holds at every $v$, the branch sequence of every orbit is a fixed
eventually-periodic sequence independent of the state; the orbit is a periodic
composition of affine maps with closed form $x_n=c\lambda^{n}+d$, and the
machine performs no state-dependent branching.
\end{theorem}

\begin{evidence}
$6{,}462$ (machine, branch) pairs, zero disagreements.  \Meas{}: in the box
$\alpha\in[0,3]$, $\beta\in[-4,8]$, $\gamma,\delta,\varepsilon\in[-3,4]$
($19{,}092$ well-defined members), exactly $2{,}351$ ($12.31\%$) are globally
rigid, all with $\beta$ even.  A closed-form criterion is exact for $\beta$
even and nonzero; the case $\beta=0$ is genuinely different
($\vp{2}(A_v)=v+1$ unbounded makes rigidity generic) and supplies $758$ of the
$2{,}351$.
\end{evidence}

\begin{remark}[the counting genre cannot work]\label{rem:counting}
$(A_v,B_v)\bmod M$ is eventually periodic in $v$ \Proved{}, and the achievable
branch tables form a rank-$5$ subgroup among tables of length $V$ \Proved{}.
It is tempting to conclude the family is too poor to host a Conway-style
reduction, which needs $M$ independent branches at modulus $M$.  That
conclusion is unsupported: measured branch \emph{bandwidth}---distinct
$(A_v,B_v)\bmod M$ actually achieved---is at least $M$ for every modulus we
tested ($100$ at $M=23$, $412$ at $M=101$, $1{,}036$ at $M=257$) \Meas{}.
Moreover the genre cannot be repaired \Interp{}: universality is a property of
a \emph{single point}, not of a family---a five-parameter family may contain a
universal member because a universal machine is a zero-parameter object---so no
dimension, rank, density or bandwidth argument can exclude it.  What survives
is a compression bound \Proved{}: at most $\lfloor\log_2(|\beta|/\rho)\rfloor$
branches are pairwise $\rho$-separated in slope, so $k$ separated branches cost
$|\beta|\ge\rho2^{k-1}$.
\end{remark}

\begin{remark}[automatic invariants]\label{rem:automatic}
\CAP{} for each stated size.  Bounded-size searches for $2$-automatic
invariants separating the orbit from the halting set return none for the
Needle's reduction at $\le11$ states (least-significant-bit-first,
minimal-word), $\le13$ (LSB, zero-invariant) and $\le13$
(most-significant-bit-first), and none for M3 in base $3$ at $\le7$ states;
encodings were shown non-vacuous by planting certificates at matched sizes.
Unbounded size is open, and after Theorem~\ref{thm:t15} this is the sharpest
open certificate question for the family.  Known non-definability results
\cite{dhiman-pandey} constrain the reachability \emph{relation} of
$3x{+}1$-type maps in B\"uchi arithmetic but do not obstruct invariants.
\end{remark}

\section{Status of every result, and threats to validity}
\label{sec:ledger}

\subsection{The ledger}

\begin{center}
\small
\begin{tabular}{@{}llp{6.4cm}@{}}
\toprule
result & status & note \\
\midrule
Prop.~\ref{prop:needle-id} identification & \Proved & one line of algebra \\
Prop.~\ref{prop:exact} exactness & \Proved & \\
Thm~\ref{thm:census} census & \CAP & each of $318$ decisions is an exact finite
  closure computation; all re-audited \\
Thm~\ref{thm:criteria} (T3), (C4) & \Proved & hand proofs given \\
Thm~\ref{thm:criteria} (C6), (C8) & \CAP & exhaustive on $(\Z_M)^5$; hand case
  analysis in repository only \\
Lem.~\ref{lem:universal} sieve lemma & \Proved & three-line reduction \\
Cor.~\ref{cor:linear} linear form & \Proved & \\
Thm~\ref{thm:ten} ten machines & \Proved & closed forms are exact identities;
  finitely many boundary branches checked \\
Thm~\ref{thm:complete} completeness & \CAP & finite: ten proofs plus $620$
  witnesses; \textbf{relative to the box} \\
Obs.~\ref{obs:nothinning} no thinning & \Meas & no asymptotic claim made \\
Thm~\ref{thm:t12} $\delta$-saturation & \Proved & \\
Lem.~\ref{lem:descent} descent & \Proved & \\
Thm~\ref{thm:t14} even saturation & \Proved & \\
Lem.~\ref{lem:reach} reach & \Proved & construction also executed \\
Thm~\ref{thm:t15}, Cor.~\ref{cor:needle} & \Proved & the paper's main result \\
Prop.~\ref{prop:m3table} M3 branch table & \MV & induction not written out;
  Thms~\ref{thm:m3n2},~\ref{thm:m3n1} depend on it \\
Lem.~\ref{lem:baseq} base-$q$ lemma & \Proved & \\
Thm~\ref{thm:m3n2} M3 two-step pinning & \CAP & least independently
  corroborated result here \\
Thm~\ref{thm:m3n1} M3 no-certificate & \Proved$^{*}$ & modulo
  Prop.~\ref{prop:m3table} \\
Rem.~\ref{rem:rank} dichotomy & \Proved & \\
Rem.~\ref{rem:boundary} boundary & \Interp & \\
Prop.~\ref{prop:sheep-id} sheep identification & \Proved & one line of algebra;
  broadly machine-checked \\
Thm~\ref{thm:sheep-sieve} sheep sieve closes & \Proved & recovers and sharpens
  a community lemma \\
Cor.~\ref{cor:sheep-H} sheep halting set & \Proved & matched to brute force
  below $3\times10^5$ \\
Thm~\ref{thm:sheep-d2} sheep depth two & \Proved & survivor set exactly
  $\{0,1,2,3,5\}$ \\
Rem.~\ref{rem:sheep-ladder} sieve saturates & \Meas & five depths; the
  convergence is extrapolated, not proved \\
Rem.~\ref{rem:sheep-dichotomy} dichotomy tested & \Interp & the prediction was
  on record first; one instance \\
Lem.~\ref{lem:t16} linear-exponential saturation & \Proved & no hypothesis on
  the multiplier \\
Thm~\ref{thm:m1n1} M1 no odd certificate & \Proved & the even-side statement
  accompanying it is \MV \\
Thm~\ref{thm:rigid} rigidity & \Proved & the $12.31\%$ figure is \Meas \\
Rem.~\ref{rem:counting} counting genre & mixed & periodicity \Proved{};
  bandwidth \Meas{}; the single-point argument \Interp \\
Rem.~\ref{rem:automatic} automatic invariants & \CAP & per stated size only \\
\S\ref{sec:discussion} placement & \Interp & a syntactic comparison, not a
  theorem \\
\bottomrule
\end{tabular}
\end{center}

\subsection{Threats to validity}
\label{sec:threats}

\begin{enumerate}
\item \textbf{All verification is by our own implementations.}  Several
  headline results were checked on two independently written code paths, and
  where possible against pre-existing implementations written for other
  purposes; but there is no third-party or formally verified check.  A
  systematic misunderstanding of the semantics would not be caught.  A Lean
  formalisation of Lemma~\ref{lem:universal}, Lemma~\ref{lem:descent} and
  Lemma~\ref{lem:reach} is the natural remedy and is not yet done.
\item \textbf{\CAP{} results depend on implementation correctness}, including
  correctness of the SAT encodings in Remark~\ref{rem:automatic}.  Encoding
  adequacy was established by planting certificates, which addresses vacuity
  but not all forms of error.
\item \textbf{Proposition~\ref{prop:m3table} is machine-verified, not proved},
  and two theorems depend on it.
\item \textbf{Theorem~\ref{thm:m3n2}'s depth-two argument received the weakest
  independent scrutiny} of any result claimed here; see its Evidence note.
\item \textbf{Selection effects in the census.}  The box was chosen for
  computational convenience.  Frequencies such as ``$29.5\%$ decided'' are
  properties of that box, not of the family.
\item \textbf{Developed against one machine---now partly mitigated.}  Every
  theorem in Sections~\ref{sec:sieve}--\ref{sec:baseq} was developed with the
  Space Needle in view, and a theory fitted to one instance can encode that
  instance's accidents.  Section~\ref{sec:wild} is the corrective: one wild
  cryptid found independently (sheep) falls inside the interface and confirms
  a prediction made beforehand, and one machine outside the family
  (M1) is nonetheless covered by the saturation argument.  Two instances is
  better than one; it is not a demonstration of generality.
\item \textbf{Asymmetry of search evidence.}  Refuting a certificate at a
  given size is cheap; finding one is expensive.  Our machinery therefore
  produces negative results far more readily than positive ones---a property
  of the evidence, not of the machines.
\end{enumerate}

\section{Progress against the goals, revisited}
\label{sec:goals-revisited}

\paragraph{G1 (observe): substantial.}
Lemma~\ref{lem:universal} and Corollary~\ref{cor:linear} give the family's
last-step sieve in closed form and isolate why it is hard---the exponential is
in the modulus, not the target.  Theorem~\ref{thm:complete} is, to our
knowledge, unusual in kind: a method closed by a completeness statement rather
than an exhausted budget.  Theorem~\ref{thm:criteria} supplies closed-form
separation criteria where only $m=3$ was previously characterised, and
Remark~\ref{rem:rank} explains a difference between two machines that had
previously only been observed.

\paragraph{G2 (decide): moved, and not on the target class.}
We decide $318+10$ members.  Every one of them is a Collatz-like function from
a manufactured box; none is a Turing machine; and each is decided precisely
because it admits a certificate, which is what it means not to be hard.  The
flagships remain open.  We regard the honest statement of this contribution as:
\emph{the census partitions a family into an easy part, a characterised hard
part, and---via Theorem~\ref{thm:t15}---a proof that the boundary between them
is not an artifact of effort}.  It should also be said plainly that
Theorem~\ref{thm:t15} makes G2 harder, not easier: it removes the cheapest
technique from consideration for the hard members permanently.

\paragraph{G3 (explain): the paper's main yield.}
Before this work the statement was ``no separating congruence was found for the
Needle's reduction at $m\le20{,}000$''---a bounded search, silent beyond its
budget.  Corollary~\ref{cor:needle} replaces it with an unconditional
impossibility, and Sections~\ref{sec:saturation}--\ref{sec:barriers} supply
mechanisms rather than only outcomes: $\delta$ sweeps every modulus; $\beta$'s
parity governs the sole $2$-adic escape hatch; rigidity marks off a slice that
cannot compute; and one genre of universality barrier is refuted with a reason
that applies to all its members.

\paragraph{G4 (transfer): demonstrated three times, with its boundary.}
Section~\ref{sec:baseq} moves the whole chain to base $3$ and obtains
Theorem~\ref{thm:m3n1} for a second machine.  Section~\ref{sec:sheep} applies
it unchanged to a cryptid found by someone else, and
Section~\ref{sec:t16} strips the saturation argument of the family schema
entirely, then applies it to a machine of a different kind
(Theorem~\ref{thm:m1n1}).  Equally informative is Remark~\ref{rem:boundary}:
two further families provably lie outside the toolkit's interface.  Three
transfers is evidence of generality, not proof of it.

\paragraph{What did not move.}
No progress was made on either flagship's actual halting question; the sheep
machine's orbit question is equally untouched, and the fact that its
\emph{arithmetic} closes completely while its orbit question does not is
itself a reminder of how little the last-step analysis reaches; no positive
non-halting theorem was obtained for any hard instance by any method; and the
universality question for the family remains open, with one route now closed.

\section{Discussion: where the difficulty now sits}
\label{sec:discussion}

Combining Lemma~\ref{lem:universal} with Theorem~\ref{thm:t15}: every
congruence-flavoured route to the Needle's reduction is closed, and the halting
condition itself is the two-generator membership~\eqref{eq:needle}.  Lifting to
the integers, a halt out of branch $v$ at value $2^{E}$ requires
$A_vk=2^{E}-B_v$ with $k\ge1$: a linear-exponential Diophantine condition in
which the unknown multiplier $k$ appears \emph{linearly} against the
exponentials $2^{v+1}$ and $2^{E}$.

\begin{observation}\label{obs:address}
\Interp{}\; Systems of linear-exponential equations over $\{2,3\}$ without such
mixed products are decidable
\cite{dong-shafrir,chistikov-mansutti-starchak}.  The system above differs from
those fragments in two syntactic respects: one product of a free variable with
an exponential, and a modulus $2^{v+1}+3$ that is not a prime power.  We offer
this as the most precise available \emph{placement} of the difficulty; it is
not a hardness theorem, and we make no claim that the difference is essential.
\end{observation}

It is natural to read the impossibility results against Skolem's conjecture
(an unsolvable purely exponential equation is unsolvable modulo some witness
\cite{bertok-hajdu}): the ten machine theorems are explicit local witnesses,
and Theorem~\ref{thm:t15} says the flagships admit no witness in the congruence
class.  We stress that this is an orientation and not an application: the
linear occurrences of $k$ and of $\delta v$ place our equations outside the
purely exponential class the conjecture concerns. \Interp{}

\subsection{Open problems}

\begin{question}
Does the Space Needle halt---equivalently, does the orbit of $6$ under
\eqref{eq:doucette} reach a power of two?  The per-branch obstruction is
exactly~\eqref{eq:needle}, and $120$ census members share its sieve group,
giving a graded family of test cases.
\end{question}

\begin{question}
Is there a $2$-automatic invariant separating the Needle's orbit from the
powers of two, at any number of states?  After Theorem~\ref{thm:t15} this is
the minimal certificate class not proved empty.
\end{question}

\begin{question}
Is the one-schema family Turing-complete---some single member with undecidable
halting set, or the parametrised problem?  Remark~\ref{rem:counting} rules out
the counting genre; the rigid slice of Theorem~\ref{thm:rigid} is provably not
where universality lives; universal generalised Collatz maps with a fixed
modulus exist \cite{kascak}, and integer piecewise-affine reachability is
delicate already in low dimension \cite{ben-amram}, but all known reductions
spend arbitrarily many independent branch coefficients where this family has
five.
\end{question}

\begin{question}
Prove that almost no start halts, with unbounded depth: for the Needle's
reduction, $\#\{x\le N \text{ halting}\}=O(N^{c})$ with $c<1$.  Depth-graded
polylogarithmic counts are known exactly for each fixed depth, and measured
branching along the backward tree peaks at $0.79<1$ \Meas{}, so a uniform
subcriticality argument is plausible.  This is the expanding-map mirror of
backward-tree counting \cite{krasikov-lagarias,kontorovich-lagarias}; note that
almost-everywhere results in the contracting regime \cite{tao} are not
transportable, the drift having the wrong sign.
\end{question}

\section*{Acknowledgments}
[To be completed.  Produced within an AI-assisted research program; see the
title-page note and Section~\ref{sec:conventions}.]

\begin{thebibliography}{CMS24}

\bibitem[BA13]{ben-amram}
A.~M. Ben-Amram.
\newblock Mortality of iterated piecewise affine functions over the integers:
  decidability and complexity.
\newblock In \emph{STACS 2013}, LIPIcs vol.~20, pages 514--525, 2013.

\bibitem[BH15]{bertok-hajdu}
C.~Bert\'ok and L.~Hajdu.
\newblock A Hasse-type principle for exponential Diophantine equations and its
  applications.
\newblock \emph{Mathematics of Computation}; arXiv:1407.6499, 2015.

\bibitem[bbcC]{bbchallenge-cryptids}
The bbchallenge collaboration.
\newblock Cryptids.
\newblock \url{https://wiki.bbchallenge.org/wiki/Cryptids}, accessed July 2026.

\bibitem[bbcS]{bbchallenge-needle}
The bbchallenge collaboration.
\newblock Space Needle.
\newblock \url{https://wiki.bbchallenge.org/wiki/Space_Needle}, accessed July
  2026.  (Turing machine due to mxdys, January 2025; one-variable reduction due
  to K.~Doucette.)

\bibitem[CMS24]{chistikov-mansutti-starchak}
D.~Chistikov, A.~Mansutti, and M.~Starchak.
\newblock Integer linear-exponential programming in NP by quantifier
  elimination.
\newblock In \emph{ICALP 2024}; arXiv:2407.07083, 2024.

\bibitem[Con72]{conway72}
J.~H. Conway.
\newblock Unpredictable iterations.
\newblock In \emph{Proceedings of the 1972 Number Theory Conference, Boulder},
  pages 49--52, 1972.

\bibitem[DP26]{dhiman-pandey}
Dhiman and Pandey.
\newblock Non-definability results for $3x{+}1$-type transition and
  reachability relations in B\"uchi arithmetic.
\newblock arXiv:2601.12772 and arXiv:2602.06066, 2026.
\newblock (Author names to be completed before submission.)

\bibitem[DS25]{dong-shafrir}
R.~Dong and D.~Shafrir.
\newblock $S$-unit equations in modules and linear-exponential Diophantine
  equations.
\newblock arXiv:2505.19141; STOC 2026, 2025.

\bibitem[Ka\v{s}92]{kascak}
F.~Ka\v{s}\v{c}\'ak.
\newblock Small universal one-state linear operator algorithm.
\newblock In \emph{MFCS 1992}, LNCS vol.~629, pages 327--335, 1992.

\bibitem[KL03]{krasikov-lagarias}
I.~Krasikov and J.~C. Lagarias.
\newblock Bounds for the $3x+1$ problem using difference inequalities.
\newblock \emph{Acta Arithmetica}, 109:237--258, 2003.

\bibitem[KL10]{kontorovich-lagarias}
A.~V. Kontorovich and J.~C. Lagarias.
\newblock Stochastic models for the $3x+1$ and $5x+1$ problems.
\newblock In \emph{The Ultimate Challenge: The $3x+1$ Problem}, AMS;
  arXiv:0910.1944, 2010.

\bibitem[KS07]{kurtz-simon}
S.~A. Kurtz and J.~Simon.
\newblock The undecidability of the generalized Collatz problem.
\newblock In \emph{TAMC 2007}, LNCS vol.~4484, pages 542--553, 2007.

\bibitem[Lag85]{lagarias-monthly}
J.~C. Lagarias.
\newblock The $3x+1$ problem and its generalizations.
\newblock \emph{American Mathematical Monthly}, 92:3--23, 1985.

\bibitem[Lag21]{lagarias-overview}
J.~C. Lagarias.
\newblock The $3x+1$ problem: an overview.
\newblock arXiv:2111.02635, 2021.

\bibitem[Tao22]{tao}
T.~Tao.
\newblock Almost all orbits of the Collatz map attain almost bounded values.
\newblock \emph{Forum of Mathematics, Pi}, 10:e12, 2022.

\end{thebibliography}

\end{document}
